\section{RAG Pipeline 实验对比}

\subsection{实验设计}

作者设计了一个定量可评估的任务——构建一个针对学术论文的 RAG（Retrieval-Augmented Generation）问答 Pipeline：

\begin{enumerate}
    \item 从 Hugging Face 每日论文中取 5 篇论文
    \item 构建包含 100 个问题和标准答案的测试集
    \item 给两个代理相同的 prompt，要求构建完整 RAG Pipeline
\end{enumerate}

给两个代理的统一要求：
\begin{itemize}
    \item Python 实现，使用 \texttt{PyMuPDF} 处理 PDF
    \item 自选分块策略（chunking strategy）
    \item 自选向量索引方案
    \item 使用 \texttt{llama-3.1-8b-instant} 生成答案
    \item 证据不足时返回 fallback 响应，不得幻觉
\end{itemize}

两个代理均使用最高配模型（Opus 4.6 / GPT-5.3-Codex），High effort 推理强度，无 \texttt{AGENTS.md}。

\subsection{实现方案对比}

\begin{table}[H]
\centering
\caption{RAG Pipeline 实现方案对比}
\begin{tabular}{lp{5cm}p{5cm}}
\toprule
\textbf{维度} & \textbf{Claude Code} & \textbf{Codex} \\
\midrule
Embedding 模型 & all-MiniLM-L6-v2 & all-MiniLM-L6-v2 \\
\addlinespace
Top-K & $k=5$ & $k=5$ \\
\addlinespace
向量存储 & ChromaDB & FAISS（更底层、内存效率更高） \\
\addlinespace
分块策略 & 递归字符分割：\texttt{\textbackslash n\textbackslash n} $\to$ \texttt{\textbackslash n} $\to$ \texttt{.} $\to$ 空格；1000 字符/200 字符重叠 & 句子级词分割，每块 220 词/40 词重叠 \\
\addlinespace
检索度量 & 原始 L2 距离 & 内积（cosine）分数 \\
\addlinespace
置信度机制 & 单阈值（L2 $> 1.2$ 视为不相关）+ 平均距离检查 & 多标准三级机制：strong / moderate / insufficient \\
\addlinespace
代码架构 & 扁平函数式，模块级常量 & OOP Pipeline 类，集中配置，dataclass，argparse CLI \\
\bottomrule
\end{tabular}
\end{table}

\begin{knowledgebox}{工程质量观察}
Codex 的实现在代码架构上明显更优：面向对象设计、集中配置管理、类型标注（dataclass）、CLI 参数解析。在大型或正式项目中，这种工程化差异非常关键。
\end{knowledgebox}

\subsection{行为差异}

\begin{itemize}
    \item \textbf{Claude Code}：自动端到端测试脚本，确保 Pipeline 开箱即用。首次运行无错误。
    \item \textbf{Codex}：完成实现后要求用户手动 \texttt{pip install} 并运行脚本。首次运行报错，经修复后正常。
    \item Codex 解释计划时更加详尽（verbose），Claude Code 更倾向于直接写代码执行命令。
    \item Codex 首个 token 的响应延迟可达 1 分钟，Claude Code 明显更快。
\end{itemize}

\subsection{评测结果}

使用 GPT-5.4 作为 LLM-as-a-Judge，在 Correctness、Completeness、Relevance、Conciseness 四个维度评估：

\begin{table}[H]
\centering
\caption{RAG Pipeline 评测结果（100 个问题）}
\begin{tabular}{lccc}
\toprule
\textbf{结果} & \textbf{Claude Code 胜出} & \textbf{Codex 胜出} & \textbf{平局} \\
\midrule
题数 & 42 & 33 & 25 \\
\bottomrule
\end{tabular}
\end{table}

\begin{importantbox}{Claude Code 胜出的关键因素}
Claude Code 胜出主要归因于：（1）更宽松的置信度阈值，避免了过多 fallback 响应；（2）稍高的生成温度（0.2 vs.\ Codex 的 0.1），增加了回答的丰富度。
\end{importantbox}

\begin{warningbox}{实验局限性}
这是一个简单的单次实验。在专业场景中，分块策略、向量数据库、检索方案等架构决策应由开发者主导，而非完全交给 AI。此实验更适合作为观察两款代理"默认行为"的参考，而非最终结论。
\end{warningbox}

\subsection{本章小结}

在 RAG Pipeline 任务中，两款代理选择了相似的基础方案但在关键细节上各有侧重。Claude Code 在端到端交付和最终评测上略优，Codex 在代码工程质量上更胜一筹。

